% budget_and_scope.tex — Budget analysis and production scope

\section{Budget and Production Scope}

\subsection{Budget Range: \filmbudgetrange}

This is a prestige drama, not a spectacle. The budget should reflect taste, not scale. The target range positions the film in the zone where serious prestige drama operates: above micro-budget, below the point where commercial return expectations distort creative decisions, and concentrated in a few recurring worlds the audience will actually remember.

\subsection{Cost Drivers}

\subsubsection{Period Production}
The film spans 1912--1924 across Prague, Berlin, Vienna, and the Czech countryside. Period production design, wardrobe, and locations are the primary cost drivers, but the dramatic architecture is contained: repeated institutions, repeated domestic interiors, and a small number of true set-piece environments.

\textbf{Mitigation}: Shoot primarily in the Czech Republic and Germany, where production infrastructure is strong, period-appropriate locations are abundant, and tax incentives are favorable.

\subsubsection{Locations}
\begin{itemize}[leftmargin=*]
\item \textbf{Prague}: Still retains significant period architecture. Barrandov Studios provides world-class soundstage capacity. Czech film incentive of 20\% cashback on qualifying expenditure.
\item \textbf{Berlin}: The Steglitz apartment scenes can be shot on a Prague soundstage or in period-appropriate Berlin locations.
\item \textbf{Austrian countryside}: The Kierling sanatorium scenes require a modest location with period character.
\item \textbf{Baltic coast}: The Müritz scenes can be achieved in northern Germany or the Baltic states.
\end{itemize}

\subsubsection{Cast}
Brody's quote will be the single largest above-the-line cost. Supporting cast can be drawn from the deep bench of European character actors, which provides extraordinary quality at manageable rates.

\subsubsection{Production Design}
The film requires detailed but contained period recreation. The primary settings --- the family apartment, the Insurance Institute, the shop, cafés, sanatoriums --- are interior-heavy and achievable through a combination of practical locations and set builds.

\subsection{Where the Money Shows}

This is a film where production value is expressed through precision, not scale:
\begin{itemize}[leftmargin=*]
\item The Kafka family apartment, the fancy goods shop, and the Insurance Institute as repeatable, high-value period worlds
\item The Berlin apartment and sanatorium rooms, where emotional intimacy has to be visible without production ostentation
\item Sound design and foley that turn paper, filing, stamping, and coughing into a system the audience can feel
\item Extended scene time and schedule protection for the elevator, engagement dinner, Letter, Brod, and final-room sequences
\end{itemize}

The spend is legible on screen because it is concentrated in atmosphere, performance support, and a handful of environments that recur until they feel mythic.

\subsection{Financing Lanes}

\subsubsection{Lane 1: Ultra-Premium Indie}
A24, MUBI, or equivalent. Full creative control, lower marketing spend, festival-driven release, strongest fit if the package closes with a top-tier auteur and sells the film as a severe actor showcase with a clear title object. Budget: \$12--14M.

\subsubsection{Lane 2: Mini-Major Prestige}
Searchlight, Focus Features. Broader theatrical release, larger marketing commitment, awards campaign infrastructure, and the best lane if the emphasis is Brody plus adult prestige accessibility. Budget: \$14--18M.

\subsubsection{Lane 3: Streamer Awards Play}
Netflix, Apple, Amazon. Larger budget tolerance, global distribution, aggressive awards spend. Budget: \$15--20M.

\subsubsection{Lane 4: International Co-Production}
Czech/German/Austrian co-production structure leveraging multiple national film funds and tax incentives. Most financially efficient model and arguably the cleanest structural fit for the material itself. Requires experienced line production and co-production legal expertise.

\subsection{Tax Incentives}

\begin{itemize}[leftmargin=*]
\item \textbf{Czech Republic}: 20\% cash rebate on qualifying Czech expenditure (one of Europe's most competitive incentives)
\item \textbf{Germany}: DFFF grants up to 20\% of German production expenditure; regional fund support available
\item \textbf{Austria}: FISA incentive of up to 30\% on Austrian production expenditure
\item \textbf{Combined potential}: A multi-territory shoot structure could reduce effective budget requirements by 15--25\%
\end{itemize}

\subsection{Revenue Model}

\begin{itemize}[leftmargin=*]
\item Domestic theatrical (platform release; awards-driven box office)
\item International theatrical (European markets; strong natural audience in German-speaking territories)
\item Prestige streaming sale (domestic and/or international)
\item International pre-sales (territory-by-territory; strong value in Europe)
\item Awards campaign cost offset through distributor partnership
\end{itemize}

Comparable reference: \textit{Tár} grossed \$26M worldwide on a \$35M budget; \textit{The Brutalist} demonstrated strong performance on a lower budget base. The financial proposition is not blockbuster returns but favorable risk-adjusted returns within the prestige segment.
